\documentclass[11pt,a4paper]{article}

% Page layout and mathematical tools
\usepackage[margin=1in]{geometry}
\usepackage{amsmath,amssymb,amsthm}
\usepackage[T1]{fontenc}
\usepackage{lmodern}
\usepackage{hyperref}
\usepackage{framed}

\newenvironment{asidebar}{%
    \def\FrameCommand{\vrule width 1.5pt \hspace{8pt}}%
    \MakeFramed{\advance\hsize-\width\FrameRestore}%
}{\endMakeFramed}
\newcommand{\aside}[1]{\begin{asidebar}\noindent\textbf{* Aside:} #1\end{asidebar}}

\title{Algebraic Curves Lecture Notes}
\author{Ryland Gross}
\date{\today}

% Number results within each section. These environments share one counter.
\theoremstyle{plain}
\newtheorem{theorem}{Theorem}[section]
\newtheorem{proposition}[theorem]{Proposition}
\newtheorem{lemma}[theorem]{Lemma}
\newtheorem{corollary}[theorem]{Corollary}
\newtheorem{exercise}[theorem]{Exercise}

\theoremstyle{definition}
\newtheorem{definition}[theorem]{Definition}
\newtheorem{example}[theorem]{Example}

\theoremstyle{remark}
\newtheorem{remark}[theorem]{Remark}

\begin{document}

\maketitle

\section{Lecture 1 (October 1st 2026)}

\noindent
The majority of the class will be working over a general field $k$ which is most often identified with $\mathbb{C}$. In fact, we will be working over characteristic $0$ fields (ie there is no such $n \in \mathbb{Z}^+$ such that $n\cdot 1 = 0$). Sometimes the course will take place over $\mathbb{R}$, but this is primarily just for visualization.

\begin{definition}
    An \emph{affine algebraic set} is the set of points on which a finite collection of polynomials $f_1, \dots f_m \in k[x_1,\dots,x_n]$ vanishes. We write this as $V_k(f_1, \dots,f_m)$. Once again, $$V_k(f_1,\dots,f_m)=\{\text{all } a \in k^n : f_i(a) = 0, \text{ for all } i\}$$ 
\end{definition}

\medskip
\noindent
We moved onto examples after this, many of which were drawn to be visualized (which I might add here later).

\begin{example}
    $V_\mathbb{R}(x^2+y^2-1)$ is precisely the unit circle since it vanishes when $x^2 + y^2 = 1$.
\end{example}

\begin{example}
    $(2,3) \in \mathbb{R}^2$, the point, is defined by $V_\mathbb{R}(x-2,y-3) = \{(2,3)\} \subseteq \mathbb{R}^2$ since these are clearly the roots of our two polynomials.
\end{example}

\begin{example}
    $\mathbb{R}^2$ is the vanishing of $0$'s constant map: $V_\mathbb{R}(0)$.
\end{example}

\begin{example}
    $\{a_1, \dots a_n\} \subseteq \mathbb{C}$, points in $\mathbb{C}$ we can write this as the vanishing of $P(x)$ with these complex numbers as roots of said polynomial. We take $(x-a_1)\dots(x-a_n) = P(x)$ and $V_\mathbb{C}(P(x))$ is now our affine algebraic set for this set in $\mathbb{C}$
\end{example}

\begin{example}
    The integers, $\mathbb{Z} \subseteq \mathbb{C}$ cannot be a affine algebraic set, because we need a finite number of finite degree polynomials which can only have finitely many roots.
\end{example}

\noindent
Now, we move on to defining a more specific curve that the class will focus on.

\begin{definition}
    An affine plane curve over $k$ is a cut-out from a single non-constant polynomial in two variables. $$C = V_k(P(x,y)) = \{(x,y) \in k^2 | P(x,y) = 0\}$$
\end{definition}

\noindent
This definition of a plane curve develops interesting questions for us to consider. \begin{enumerate}
    \item How can we descrive the set of points of a given plane curve?
    \item What can we say about the intersection of two plane curves? (Bezout's Theorem)
    \item When do two polynomials describe the same plane curve? (Nullstellensatz)
\end{enumerate}

\noindent
Let's look at that last question carefully. We got some examples in class about this, and I will expand on it a little bit. In $\mathbb{R}$, we have $V_\mathbb{R}(x^2+1) = \emptyset$, and we have $V_\mathbb{R}(1) = \emptyset$ as well, but they have very little in common as polynomials. The interesting fact, however, is that this happens precisely because $\mathbb{R}$ is not algebraically closed, over $\mathbb{C}$ this is \emph{less} of a problem. When we examine $V_\mathbb{C}(x^2+1) = \{-i,i\}$ and $ V_\mathbb{C}(1) = \emptyset$ we see clearly that we have gotten rid of the problem. 

\medskip
\noindent
We then define the degree of a polynomial which is the largest such $i+j$ you can find in the exponents of the polynomial $$P(x,y) = \sum_{i,j}a_{ij}x^iy^j$$ The different kinds of plane curves we have are as follows \begin{enumerate}
    \item \textbf{Degree 1:} These are linear equations describable by $ax+by+c$ and the vanishing of that polynomial where $a, b$ not both $0$
    \item \textbf{Degree 2:} These, like in high-school, are called conics and can be described generally by the vanishing of $P(x,y) = ax^2+bxy+cy^2+dx+ey+f$. A classic example is the unit circle, the vanishing of $x^2+y^2-1$. 
    \item \textbf{Degree 2: Ellipse:} for example, the vanishing of $P(x,y) = x^2+4y^2-1$.
    \item \textbf{Degree 2: Hyperbolas:} for example, the vanishing of $P(x,y) = x^2-y^2-1$.
    \item \textbf{Degree 2: Parabolas:} for example, the vanishing of $P(x,y) = x^2-y$ 
    \item \textbf{Degree 2: Degenerate Examples:} for example, the vanishing of $xy$ or $x^2$ or $y^2$ which respecively, vanish on the union of the $x$ and $y$ axes, $x$ axis, and $y$ axis. 
\end{enumerate}

\noindent
Now, we might be interested in certain solutions to intersections of lines with some of our plane curves. We look at an example from antiquity related to Pythagorean triples. We start by turning the Pythagorean theorem into a plane curve, $$a^2 +b^2 = c^2 \implies \frac{a^2}{c^2} + \frac{b^2}{c^2} = 1$$ Hence, if we have pick $x = a/c$ and $y = b/c$ we can have $x^2+y^2=1$ be our plane curve. So, we're really looking to rational points on the unit circle. Thus, pick $p_0 = (-1,0)$ and we will examine a line going through the unit circle as a chord (secant line) and we will examine the solution to the second intersection which gives us a parameterization for the rest of the circle. This curve will look like $y=t(x+1)$. Now, we have two equations $$x^2+y^2=1 \qquad y=t(x+1)$$ Subsitution helps us get $$x^2+t^2(x+1)^2=1 \implies (x+1)(x-1+t^2(x+1))=0 \implies x = \frac{1-t^2}{1+t^2}$$ Now, solving for $y$ once again with subtitution we get, $$y = \frac{2t}{1+t^2} \implies \text{Unit-Circle} = \left\{(-1,0) \cup \left(\frac{1-t^2}{1+t^2},\frac{2t}{1+t^2}\right), t \in \mathbb{Q}\right\}$$ Since we want rational $t$ we can pick $t= r/s$ and get, $$\text{Unit-Circle} = \left\{(-1,0) \cup \left(\frac{s^2-r^2}{s^2+r^2},\frac{2rs}{s^2+r^2}\right)\right\}$$ Which means we know our $a$ and $b$ for checking pythagorean triples which is, (after multiplying over by $c$ to clear denominators) $$(s^2-r^2, 2rs, s^2+r^2), r,s \in \mathbb{Z}$$ But it is important to note that we will need to pick $s, r$ with different parity or else we will have a triple that is all even. For example see if we pick $5, 7$ for our $s$ and $r$ we get $$(7^2-5^2, 2\cdot 5 \cdot 7, 7^2+5^2) = (24, 70, 74)$$ which we can factor and see $12, 35, 37$ is the triple that actually defines this one. Thus, what we need is not merely two co-prime integers, we want $\gcd(s,r)=1$ and $s \mod 2 \neq r \mod 2$. Choosing these will give you infinitely many unique pairs since for any $a$ which is even, there are infinitely many prime numbers $p$ that are larger than $a$ and hence coprime. 

\section{Lecture 2 (October 2nd 2026)}

The majority of this class was reviewing some ideas from algebra and topology. We discuss ideals, the Hilbert Basis Theorem. Afterwards, we look at some algebra we can perform over varieties.

\begin{definition}
    An \emph{ideal} in a commutative ring $R$ is $I \subseteq R$ which has two properties \begin{enumerate}
        \item Closure under addition: for any $a,b \in I$ we have $a+b \in I$
        \item Closure under multiplication by elements of $R$ such that for any $r \in R, a \in I$ we have $ra \in I$
    \end{enumerate}
\end{definition}

\begin{exercise}
    The Ideal generated by a subset $X\subseteq R$ is the smallest ideal that contains all of $X$.
\end{exercise}

\begin{proof}
    Placeholder
\end{proof}

\begin{example}
    The ideal generated by $(1)$ the unit of $R$ is equal to all of $R$.
\end{example}

\begin{example}
    $R = \mathbb{Z}$ is a ring and we have ideals $(n)$ for $n \in \mathbb{Z}$ because $\mathbb{Z}$ is a euclidean domain and hence a PID.
\end{example}

\begin{example}
    in the ring $R = k[x,y]$ we have $(x-3)$ and $(x-3,y-2)$ as ideals.
\end{example}

\aside{ The original hilbert basis theorem is not proven in our class, so we introduce and prove that here for completion's sake. 
\begin{definition}
    A \emph{Noetherian} ring $R$ is a ring that satisfies the ascending chain condition: for any sequence of Ideals (no left/right since we are only worried about commutative rings) $I_1 \subseteq I_2 \subseteq \dots \subseteq I_n \subseteq \dots$ there exists $m \in \mathbb{N}$ such that $I_m = I_{m+1} = I_{m+2} = \dots$. We say that the chain stabilizes. 
\end{definition}

\begin{theorem}[Hilbert Basis Theorem]
    If $R$ is a noetherian ring then $R[X]$ is noetherian ring.
\end{theorem}

\begin{proof}
    Placeholder
\end{proof}}

\begin{theorem}[Our Class' Hilbert Basis Theorem]
    For any field $k$, any ideal $I \subseteq k[x_1, \dots x_n]$ has a finite generating set such that $I = (f_1, f_2, \dots f_l)$
\end{theorem}

\aside{This version was also not proven in class so we will do it here: \begin{proof}
    Placeholder
\end{proof}}

\noindent
If $I \subseteq k[x_1, \dots, x_n]$ is an ideal then we write $$V_k(I) = \{(a_1, \dots a_n) \in k^n | f(a_1, \dots a_n) = 0, \text{for all} f \in I\} $$ Additionally by the Hilbert Basis Theorem since we know that $I$ is finitely generated we have $$I = (f_1, \dots, f_l)$$ We can also check that $$I = \{g_1f_1+g_2f_2 +\dots + g_lf_l : \text{for any} g_1, \dots, g_l \in k[x_1, \dots,x_n]\}$$ This helps motivate our next observation.

\begin{remark}
    If $I = (f_1, \dots f_l)$ then $V_k(I) = V_k(f_1, \dots, f_l)$ 
\end{remark}

\begin{proof}
    Placeholder
\end{proof}

\begin{proposition}
    Let $V_1$ and $V_2$ be affine algebraic sets in $k^n$ then $V_1 \cap V_2$ is an affine algebraic set.
\end{proposition}

\begin{proof}
    Placeholder
\end{proof}

\begin{exercise}
    Let $V_1$ and $V_2$ be affine algebraic sets in $k^n$ then $V_1 \cup V_2$ is an affine algebraic set.
\end{exercise}

\begin{proof}
    Placeholder
\end{proof}

\noindent
Quick check! Is the intersection of two affine plane curves also an affine plane curve? No. This is a key distinction between plane curve and affine algebraic set. The key distinction is that an affine algebraic set is defined by the vanishing of finitely many polynomial functions $f_1, \dots f_m \in k[x_1, \dots x_n]$ but a plane curve is the vanishing of a single non-constant polynomial in two variables. Thus, we can look at $$P(x,y) =x \quad Q(x,y) = y \implies V_k(P) \cap V_k(Q) = (0,0) $$ Thus, for this to be a plane curve it has to be the vanishing of a single non-constant polynomial which is impossible. 

\begin{remark}
    The intersection over some indexing set $\bigcap_\alpha V(I_\alpha)$ of arbitrarily many affine algebraic sets is still an affine algebraic set because $I_\alpha = (f_{\alpha,1}, \dots f_{\alpha, i})$ so you can just look at the vanishing of all the generators $$ \bigcap_\alpha V(I_\alpha) = V\left(\sum_{\alpha} I_\alpha\right) \text{ since we have } \sum_\alpha I_\alpha = I\left(\bigcup_\alpha X_\alpha\right) \text{ is finitely generated by 2.8}$$ 
\end{remark}

\noindent
This remark lets us define the Zariski topology on affine space. Here we define the open sets on affine space as the complements of affine algebraic sets. Let's recall what we need to have for a topology, \begin{enumerate}
    \item $\emptyset$ and $k^n$ are open sets.
    \item finite intersections of open sets are open
    \item arbitrary unions of open sets are open
\end{enumerate}

\begin{exercise}
    Prove the \emph{Zariski Topology} on affine space is indeed a topology.
\end{exercise}

\begin{proof}
    Placeholder
\end{proof}

\noindent
This topology is weaker (she actually means coarser) than the typical topologies on $\mathbb{R}^n$ and $\mathbb{C}^n$. Indeed Zariski opens are open in the typical sense but open sets in the typical topologies are not necessarily open in Zariski topology. 

\aside{\begin{definition}
    A topology on $X$ is called Hausdorff if for any two points $a,b \in X$ there exist $U, V$ open sets such that $a \in U$ and $b\in V$ and $U \cap V = \emptyset$ 
\end{definition}}

\begin{exercise}
    The Zariski Topology is not Hausdorff.
\end{exercise}

\begin{proof}
    Placeholder
\end{proof}

\begin{theorem}[Fundamental Theorem of Algebra]
    $\mathbb{C}$ is algebraically closed.
\end{theorem}
\aside{
\begin{proof}
    Placeholder (may remove)
\end{proof}
}
\begin{proposition}
    if $k$ is algebraically closed then any plane curve $C \subseteq k^2$ has infinitely many points over $k$; if $k$ is $\mathbb{C}$ then the curve is unbounded. 
\end{proposition}

\section*{Lecture 3}

\active 

\end{document}
